% 齿轮传动设计准则流程图
\begin{tikzpicture}[
    node distance=1.2cm and 1.5cm,
    >=latex,
    box/.style={
        rectangle,
        draw,
        fill=#1,
        minimum height=0.8cm,
        minimum width=3.2cm,
        align=center,
        font=\small\bfseries
    },
    box/.default=blue!20,
    graybox/.style={
        rectangle,
        draw,
        fill=gray!40,
        minimum height=0.7cm,
        minimum width=3cm,
        align=center,
        font=\small
    },
    greenbox/.style={
        rectangle,
        draw,
        fill=green!15,
        minimum height=0.7cm,
        minimum width=3.5cm,
        align=center,
        font=\small
    },
    arrow/.style={
        ->,
        thick,
        >=latex
    }
]

% 顶部节点
\node[box] (top) {闭式齿轮传动};

% 第二层节点
\node[box, below left=1.2cm and 1.5cm of top] (soft) {软齿面齿轮\\(HBS$\leqslant$350)};
\node[box, below right=1.2cm and 1.5cm of top] (hard) {硬齿面齿轮\\(HBS$>$350)};

% 第三层节点
\node[graybox, below=of soft] (soft_fail) {齿面疲劳点蚀};
\node[graybox, below=of hard] (hard_fail) {齿根弯曲疲劳折断};

% 第四层节点
\node[greenbox, below=of soft_fail] (soft_rule) {齿面接触疲劳强度准则};
\node[greenbox, below=of hard_fail] (hard_rule) {齿根弯曲疲劳强度准则};

% 开式齿轮传动
\node[box, below left=1.2cm and 2cm of soft_rule] (open) {开式齿轮传动};
\node[graybox, right=of open] (open_fail) {齿面磨损};
\node[greenbox, below right=0.8cm and 0.5cm of open_fail] (open_rule) {增大m和降低许用弯曲应力};

% 连接线
\draw[arrow] (top) -- (soft);
\draw[arrow] (top) -- (hard);
\draw[arrow] (soft) -- (soft_fail);
\draw[arrow] (hard) -- (hard_fail);
\draw[arrow] (soft_fail) -- (soft_rule);
\draw[arrow] (hard_fail) -- (hard_rule);
\draw[arrow] (open) -- (open_fail);
\draw[arrow] (open_fail) -- (open_rule);
\draw[arrow] (open_fail) -- (hard_rule);

\end{tikzpicture}
